\chapter{13}\label{section-12}

«Ey, still, sei still, psst,» eine weibliche Stimme hinter dem grellen
Licht flüsterte, «bin nur ich, alles ok,» und damit sah Louis wie die
Taschenlampe sich drehte und ein junges, weibliches Gesicht beleuchtete.
Er setzte sich ruckartig auf und starrte die Frau an. Sie legte ihre
linke Hand über ihre Brust.

«Bin Jess, und du? Schiebst Platte?»

Sie lächelte wie jemand, der Besuch ins Haus bittet. Ihre Stimme war
auffallend hell. Umrahmt von kurzen, zerzausten Haaren schauten Louis
übergrosse, wasserblaue Augen an. Ihr mageres Gesicht sah bleich aus.
Louis war sie unheimlich.

«Ehm,» er räusperte sich und wieder fehlte ihm ein Name, «Tim, hallo.
Platte?»

Er schaute Jess unsicher an.

«Ach, ist doch so, hast ’nen Schlafplatz gebraucht, oder?»

Sie runzelte die Stirn und schmunzelte wissend. Louis versuchte, ihren
Blick zu deuten. Sah er ein Flehen darin? Oder Angst, Neugierde, oder
alles auf einmal?

Jetzt kniff sie die Augen zusammen.

«Oder was sonst? Was machst du hier?»

Einmal mehr glaubte Louis, er würde fremdgesteuert. Eine neue Geschichte
breitete sich vor ihm aus. Er brauchte sie nur noch in passende Worte zu
kleiden.

«Ich suche nach neuen Erfahrungen,» er räusperte sich, «für vier Wochen,
vorerst. Ich will mal aus der Komfortzone ausbrechen. Um andere Seiten
des Lebens zu sehen. Weg von Sicherheit und diesem geordneten,
langweiligen Alltag. Menschen sehen, erfahren, wie andere leben und so,
es ist ’ne Art Studie.»

Louis beobachtete Jess genau, während er erzählte. Mit offenem Mund
hörte sie ihm zu. Sie sass regungslos und perplex vor ihm. Als käme er
von einem anderen Stern. Schliesslich klaubte sie ein Päckchen
Zigaretten aus ihren Shorts. Sie streckte es Louis entgegen. Er sah ihre
langen, hellen Finger.

«Nimm nur.»

Sie schaute wie ein kleines Mädchen.

Louis zögerte. Dann nahm er eine Zigarette aus der Packung. Jess gab ihm
Feuer.

Viele Jahre hatte er nicht mehr geraucht. Es tat gut. Und auch nicht.

Und dann sass \emph{er} stumm da und hörte \emph{ihr} zu. Und zog
zwischendurch an dieser Zigarette.

Es war unerträglich. Was, und wie sie erzählte. Und je länger die
Taschenlampe ihr Gesicht beleuchtete, desto klarer entdeckte er Spuren
darin. Spuren von Zerstörung und Leid. Wie dünn und zerbrechlich sie
aussah.

Bereits die dritte Nacht nutze sie den Rohbau als Nachtlager, erklärte sie.

Louis hatte ihre Anwesenheit nicht bemerkt. Jess musste gelernt haben,
sich unsichtbar zu machen.

Die Worte drängten mit einem Mal wie ein Wasserfall aus ihr heraus. Als
würde sie gerade von einem jahrelangen Redeverbot entfesselt. Jess’
atemlos geflüsterte Sätze, aneinandergereiht als ihre Geschichte, trafen
Louis.

Ein junges Leben in Verlorenheit. Angefangen mit einer kranken Mutter,
einem abwesenden Vater, Aufenthalten in Pflegefamilien, einer
mehrmonatigen Verstummung als Fünfjährige und schliesslich, seit Wochen,
als Untergetauchte.

Hätte sie nicht vor ihm gesessen, er hätte ihre Vergangenheit als
klischeehafte Aneinanderreihung von Widrigkeiten abgetan. Es
erschütterte ihn, dass diese Sicht hier nicht griff. Jess wirkte
glaubhaft.

«Ich suchte mir, was ich brauchte, bekam viel geschenkt,» dabei legte
sie ihre Handflächen zusammen. Die dazwischen eingeklemmte Taschenlampe
zündete flackernd an die Decke.

Louis erfuhr, dass Jess nach traumatischen Erlebnissen in der
Prostitution eines Tages von einem Streetworker angesprochen worden sei.
Er habe ihr von «NGO Heartwings» in Zürich erzählt und sie aufgemuntert,
sich bei der Organisation zu melden. Die Hilfe existiere schon seit
einigen Jahren. Das Gründerpaar hätte ein Programm entwickelt, um
willige Frauen beim Ausstieg aus der Szene zu unterstützen. Auch ihr
bester Freund habe ihr damals geraten, sich damit eine Chance zu geben.

Sie redete immer schneller.

Jess hatte es gewagt, sagte sie. Den Mann damals vor dem «Restaurant
Sonne» an der Langstrasse getroffen. Ungefährlich sei der gewesen,
freundlich.

«Der hat mir Kleider gegeben, Secondhandsachen, was soll’s, schau, wie
schön, ist auch von ihm, war mal verdammt teuer,» schwelgte sie und
deutete auf ihr schmutziges Shirt mit dem roten Gucci-Schriftzug. Louis
bemühte sich, in der Dunkelheit etwas zu sehen.

«Der kam dann mit Arbeitsintegration und so’m Zeug, fragte nach meinen
Kinder- und Jugendträumen. Hä, dachte, der spinnt jetzt, ging nicht mehr
hin, hatte Panik vor diesen Leuten, die mir angeblich helfen sollten.
Ich galt in dem Programm irgendwie als gescheitert, sagte mir die Frau,
scheissegal, nichts wie weg, dachte ich, mach dich dünn.»

Jess wirkte geschafft. Sie blickte abwesend in den dunklen Raum.

Ihr Zutrauen verwirrte ihn. Entweder war sie bedürftiger als er sich
vorstellen konnte, oder aber sie sah in ihm als Mann keine Gefahr.
Beides empfand Louis als Bedrohung. Als diffuse Aufforderung, irgend
etwas zu tun.

«Ja, und dann, dann ist’s halt passiert, willst du nicht wissen, Tim,»
Jess winkte ab, «die haben mich erwischt, dumm gelaufen, musste ganz
untertauchen.»

Jess holte Luft. Louis auch.

«Aber, ich hab Glück,» ihre Lippen formten ein schiefes Lächeln, «hat
mich hier keiner gefunden. Ausser dir.»

Und Jess rückte näher an ihn heran.
